\documentclass[a4paper,11pt]{article}
\def\email#1{{\tt#1}}

\def\N{\mbox{\rm{I}\hspace{-.15em}\rm{N}}}
\def\Z{\mbox{\rm{Z}\hspace{-.28em}\rm{Z}}}
\def\RR{\mbox{\rm{I}\hspace{-.15em}\rm{R}}}
\def\C{\mbox{\rm{l}\hspace{-.50em}\rm{C}}}
\def\Q{\mbox{\rm{l}\hspace{-.50em}\rm{Q}}}
\def\F{\mbox{\rm{I}\hspace{-.20em}\rm{F}}}

\usepackage{amssymb,amsmath}
\usepackage{url}
\usepackage{fullpage}
\begin{document}

\newtheorem{theo}{Theorem}
\newtheorem{coro}{Corollary}[theo]
\newtheorem{lemme}{Lemma}
\renewcommand{\thecoro}{\thetheo.\arabic{coro}}



\parindent=0cm


\section*{\center IN100 Initiation au Python\\
Contrôle continu 2}

\vspace{1 cm}

\begin{table}[h!]
\centering
\renewcommand{\arraystretch}{1.4} % augmente la hauteur des lignes
\begin{tabular}{|p{10cm}|p{3cm}|p{2cm}|}
\hline
\textbf{Étudiant :} & \textbf{Heure Début }  & \textbf{Note (/5) }  \\
& & \\
\hline
\textbf{Correcteur :} & \textbf{Heure Fin } & \\
& & \\
\hline

\hline
\end{tabular}
\end{table}

\vspace{-0.5 cm}

\subsection*{Consignes:} Vous avez 20 minutes pour coder la solution. Seules les types et les instructions vues en cours sont autorisées. Tout usage d'une fonction 
venant d'une bibliothèque qui n'a pas été vue en cours ne sera pris en compte. 

\vspace{-0.3 cm}
\subsection*{Sujet:}

Un billet de loterie est donné par 6 nombres entiers compris entre 1 et 48. 


\begin{itemize}
\item[1.]  En utilisant les notions vues en cours, donnez une fonction permettant de générer un billet quelconque pour un joueur. 
Nous rappelons que la fonction \texttt{random.randint(a,b)} permet de générer des nombres aléatoires entre \texttt{a} et \texttt{b}. 

Exemple d'usage: 
\begin{verbatim}
import random

print(random.randint(3,9))

\end{verbatim}

\item[2.] On fixe le billet  \texttt{B=[3,34,23,11,42,29]}.  Un billet est gagnant s'il contient au moins trois nombres de \texttt{B}. 
\'Ecrire une fonction permettant de vérifier, pour un billet quelconque, si celui-ci est gagnant. Si le billet est gagnant, votre fonction doit renvoyer tous les nombres 
que ce billet a en commun avec le billet \texttt{B}. Sinon, elle renverra \texttt{None}. 
\end{itemize}



\begin{table}[h!]
\centering
\renewcommand{\arraystretch}{1.4} % augmente la hauteur des lignes
\begin{tabular}{|p{15cm}|}
\hline
\textbf{Remarques étudiant :}    \\
 \\
\hline
\textbf{Remarques correcteur :}  \\
 \\
 \hline
\end{tabular}
\end{table}

\newpage


\section*{\center IN100 Initiation au Python\\
Contrôle continu 2}


\begin{table}[h!]
\centering
\renewcommand{\arraystretch}{1.4} % augmente la hauteur des lignes
\begin{tabular}{|p{10cm}|p{3cm}|p{2cm}|}
\hline
\textbf{Étudiant :} & \textbf{Heure Début }  & \textbf{Note (/5) }  \\
& & \\
\hline
\textbf{Correcteur :} & \textbf{Heure Fin } & \\
& & \\
\hline

\hline
\end{tabular}
\end{table}

\subsection*{Consignes:} Vous avez 20 minutes pour coder la solution. Seules les types et les instructions vues en cours sont autorisées. Tout usage d'une fonction 
venant d'une bibliothèque qui n'a pas été vue en cours ne sera pris en compte. 


\subsection*{Sujet:}

\begin{itemize}
\item[1.] En utilisant une seule ligne de code, créez une liste contenant tous les entiers compris entre 1 et 7. 
\begin{verbatim}
liste_initiale=[1,2,3,4,5,6,7]
\end{verbatim}

\item[2.] \'Ecrire une fonction qui prend en entrée une liste de taille quelconque, décale toutes les éléments de la liste de $i$ positions vers la gauche, de manière circulaire.   Ainsi premières $i$ éléments se 
retrouvent à la fin de la liste. 

Par exemple, pour la liste crée à la question 1 un décalage de 3 donnera la liste suivante:

\begin{verbatim}
liste_decalee=[4,5,6,7,1,2,3]
\end{verbatim}

\item[3.] \'Ecrire une fonctions qui prend en entrée une matrice de $m$ lignes et $n$ colonnes et modifie cette matrice comme suit:  la première ligne est décalée de manière circulaire de 1, la deuxième ligne est décalée de manière circulaire de 2, ..., 
l'avant dernière ligne est décalée circulairement de $m-1$ et la dernière ligne est inversée. 
La fonction devra renvoyer  \texttt{None}. 
\end{itemize}

\begin{table}[h!]
\centering
\renewcommand{\arraystretch}{1.4} % augmente la hauteur des lignes
\begin{tabular}{|p{15cm}|}
\hline
\textbf{Remarques étudiant :}    \\
 \\
 \\
\hline
\textbf{Remarques correcteur :}  \\
 \\
 \\
\hline
\end{tabular}
\end{table}


\newpage


\section*{\center IN100 Initiation au Python\\
Contrôle continu 2}


\begin{table}[h!]
\centering
\renewcommand{\arraystretch}{1.4} % augmente la hauteur des lignes
\begin{tabular}{|p{10cm}|p{3cm}|p{2cm}|}
\hline
\textbf{Étudiant :} & \textbf{Heure Début }  & \textbf{Note (/5) }  \\
& & \\
\hline
\textbf{Correcteur :} & \textbf{Heure Fin } & \\
& & \\
\hline

\hline
\end{tabular}
\end{table}

\vspace{-0.4cm}
\paragraph*{Consignes:} Vous avez 20 minutes pour coder la solution. Seules les types et les instructions vues en cours sont autorisées. Tout usage d'une fonction 
venant d'une bibliothèque qui n'a pas été vue en cours ne sera pris en compte. 

\vspace{-0.2cm}
\paragraph*{Sujet:} On lance deux dés à six faces et on note leurs résultats $a,b$.  Les différentes cas que l'on peut obtenir sont:
\begin{itemize}
\item  PAIRE: deux dés identiques et un autre différent: ex 3,3 ou 4,4
\item  SUPER : une des deux valeurs est 6 (ex 2,6)
\item SUPER EXTRAORDINAIRE les deux valeurs sont 6 (ex 6,6) 
\item NADA: dans tous les autres cas. 
\end{itemize}

Les premières trois cas sont gagnants. Le but est d'écrire un programme qui simule le comportement décrit précédemment.
\begin{itemize}
\item[1.] Créer une fonction \texttt{est\_gagnant} qui prend en paramètre un tuple de 2 valeurs et renvoie \texttt{True} si le tuple est gagnant, \texttt{False} sinon. 
\item[2.]  \'Ecrire une fonction qui génére plusieurs tuples de nombres aléatoires entre 1 et 6, en utilisant la fonction \texttt{random.randint}, et qui s'arrête lorsque un tuple gagnant est trouvé. 
La fonction renverra le message associé au tuple gagnant et le nombre de fois les dès ont été "jouées". Elle utilisera la fonction de la question 1. 
\end{itemize}

Exemple d’utilisation :

\begin{verbatim}
import random
print(random.randint(2, 5)) # affichera une valeur entre 2 et 5 inclus
\end{verbatim}


%\vspace{-0.5 cm}
\begin{table}[h!]
\centering
\renewcommand{\arraystretch}{1.4} % augmente la hauteur des lignes
\begin{tabular}{|p{15cm}|}
\hline
\textbf{Remarques étudiant :}    \\
 \\
 \\
\hline
\textbf{Remarques correcteur :}  \\
 \\
\hline
\end{tabular}
\end{table}


\newpage

\section*{\center IN100 Initiation au Python\\
Contrôle continu 2}


\begin{table}[h!]
\centering
\renewcommand{\arraystretch}{1.4} % augmente la hauteur des lignes
\begin{tabular}{|p{10cm}|p{3cm}|p{2cm}|}
\hline
\textbf{Étudiant :} & \textbf{Heure Début }  & \textbf{Note (/5) }  \\
& & \\
\hline
\textbf{Correcteur :} & \textbf{Heure Fin } & \\
& & \\
\hline

\hline
\end{tabular}
\end{table}

\subsection*{Consignes:} Vous avez 20 minutes pour coder la solution. Seules les types et les instructions vues en cours sont autorisées. Tout usage d'une fonction 
venant d'une bibliothèque qui n'a pas été vue en cours ne sera pris en compte. 



\subsection*{Sujet:}
Afin de s'entraîner pour une compétition sportive, Arthur doit pratiquer plusieurs activités sportives par jour, pendant $n$ jours. On utilise une liste imbriquée pour représenter 
les activités qu'il pratique chaque jour. Par exemple pour $n = 5$ jours, on
a:  

\begin{verbatim}
planning_activités=[["jogging", "yoga", "natation"], [], ["natation", "tennis"], 
["tennis"], ["natation", "badminton"]].
\end{verbatim}

Cela veut dire que le jour 1 il a pratiqué le jogging, le yoga et la natation, le jour 2 il n'a pratiqué aucune activité etc. 
\begin{itemize}
\item[1.] Écrire une fonction \texttt{nombre\_jours\_activites} qui prend en argument le planning et une activité et renvoie
le nombre de jours pendant lesquels cette activité a été pratiquée.
\item[2.]  Écrire une fonction \texttt{total\_activites} qui prend en argument le planning et renvoie
le nombre total d'activités pratiquées sur toute la période. 
\end{itemize}


\vspace{1 cm}

\begin{table}[h!]
\centering
\renewcommand{\arraystretch}{1.4} % augmente la hauteur des lignes
\begin{tabular}{|p{15cm}|}
\hline
\textbf{Remarques étudiant :}    \\
 \\
 \\
\hline
\textbf{Remarques correcteur :}  \\
 \\
 \\
\hline
\end{tabular}
\end{table}

\newpage

\section*{\center IN100 Initiation au Python\\
Contrôle continu 2}


\begin{table}[h!]
\centering
\renewcommand{\arraystretch}{1.4} % augmente la hauteur des lignes
\begin{tabular}{|p{10cm}|p{3cm}|p{2cm}|}
\hline
\textbf{Étudiant :} & \textbf{Heure Début }  & \textbf{Note (/5) }  \\
& & \\
\hline
\textbf{Correcteur :} & \textbf{Heure Fin } & \\
& & \\
\hline

\hline
\end{tabular}
\end{table}
\vspace{-0.3 cm}
\subsection*{Consignes:} Vous avez 20 minutes pour coder la solution. Seules les types et les instructions vues en cours sont autorisées. Tout usage d'une fonction 
venant d'une bibliothèque qui n'a pas été vue en cours ne sera pris en compte. 


\subsection*{Sujet}

Une file est une structure de donnée basé sur le principe que les premiers éléments ajoutés dans la file seront les premiers à en être retirés.
Dans cet exercice, nous utiliserons une liste pour représenter une file pour gérer les patients dans une clinique. Chaque patient est représenté par un entier correspondant au temps estimé de consultation en minutes.

\begin{itemize}
\item[1.] Écrire une fonction qui prend en argument une file et un entier, et qui rajoute cet élément dans la file, donc en début de liste. 
Vous utiliserez cette fonction pour rajouter dans votre file les temps de consultation suivants:
\begin{verbatim}
15
20
15
30
25
\end{verbatim}
\item[2.]  Écrire une fonction qui prend en argument la liste donnant la file et qui enlève le dernier élément de la liste et renvoie cet élément. 
\item[3.]  Écrire une fonction \texttt{traiter\_patients} prend en argument  la file de patients, traite tous les patients dans la file jusqu’à ce qu’elle soit vide et 
renvoie le temps total passé à consulter tous les patients.
\end{itemize}



\vspace*{-0.2 cm}
\begin{table}[h!]
\centering
\renewcommand{\arraystretch}{1.4} % augmente la hauteur des lignes
\begin{tabular}{|p{15cm}|}
\hline
\textbf{Remarques étudiant :}    \\
 \\
\hline
\textbf{Remarques correcteur :}  \\
 \\
\hline
\end{tabular}
\end{table}





\end{document}
